\documentclass[twoside, 10pt, a4paper]{article}
\usepackage[utf8]{inputenc}
\usepackage[T1]{fontenc}
\usepackage[portuguese]{babel}

% --- FONTE PADRÃO (Estilo Didático Encorpado) ---
\usepackage{helvet} 
\renewcommand{\familydefault}{\sfdefault}

% --- GEOMETRIA E ESPAÇAMENTO ---
\usepackage[top=1.6cm, bottom=1.5cm, left=0.9cm, right=0.9cm, headheight=22pt, headsep=0.4cm]{geometry}
\usepackage{microtype} 
\usepackage{setspace}
\setstretch{1.0}
\usepackage{parskip}
\setlength{\parskip}{2pt}
\setlength{\parindent}{0pt}

\newlength{\espacoPosTitulo}\setlength{\espacoPosTitulo}{0.3cm}
\newlength{\espacoPreQuestao}\setlength{\espacoPreQuestao}{0.1cm}
\newlength{\espacoQuestao}\setlength{\espacoQuestao}{0.2cm}
\newlength{\espacoAlt}\setlength{\espacoAlt}{0.2cm}

\usepackage{enumitem}
\setlist{itemsep=0.4cm, topsep=0.1cm, parsep=0pt, partopsep=0pt}
\newcommand{\larguraTikz}{0.85\linewidth} 

% --- MATEMÁTICA E CORES ---
\usepackage{amsmath, amsfonts, amssymb, nccmath, mathastext}
\usepackage{xcolor}
\definecolor{indigo}{HTML}{4F46E5}
\definecolor{CorTitUm}{HTML}{FF6F61}
\definecolor{CorTitDois}{HTML}{FF6F61}
\definecolor{CorTitTres}{HTML}{658894}
\definecolor{CorLinha}{HTML}{B0B0B0} 
\definecolor{metal}{HTML}{7F8C8D}
\definecolor{vidro}{HTML}{3498DB}
\definecolor{ouro}{HTML}{F1C40F}
\definecolor{concreto}{HTML}{95A5A6}
\definecolor{madeira}{HTML}{D35400}
\definecolor{CinzaEscuro}{HTML}{4A4A4A} 
\definecolor{CinzaMedio}{HTML}{848484} 
\definecolor{Cinzablack}{HTML}{353535} 

% --- MINI-CABEÇALHO E RODAPÉ AUTORAL (TWOSIDE) ---
\usepackage{fancyhdr}
\pagestyle{fancy}
\fancyhf{} 
\renewcommand{\headrulewidth}{0.3pt} 
\renewcommand{\footrulewidth}{0.3pt}
\fancyhead[LE]{\small\sffamily\bfseries\color{gray!75} DECIMAIS, POTENCIAÇÃO E RADICIAÇÃO}
\fancyhead[RO]{\small\sffamily\color{gray!75} \texttt{profraphaelpaulino.com.br}}
\fancyfoot[LE]{\small \textbf{\thepage}}
\fancyfoot[RO]{\small \textbf{\thepage}}

% --- CAIXAS DE DESTAQUE ---
\usepackage[most]{tcolorbox}
\newtcolorbox{boxresponda}{colback=white, colframe=gray!45, boxrule=0pt, leftrule=1.7pt, sharp corners, enlarge left by=0.4cm, width=\linewidth-0.4cm, left=8pt, right=0pt, top=2pt, bottom=2pt, fontupper=\small}
\definecolor{CorDicaBorda}{HTML}{17c1be} 
\definecolor{CorDicaFundo}{HTML}{f3fbfb} 
\newtcolorbox{boxdica}{colback=CorDicaFundo, colframe=CorDicaBorda, boxrule=0pt, leftrule=2.5pt, sharp corners, width=\linewidth, left=8pt, right=8pt, top=6pt, bottom=6pt, halign=center, fontupper=\small}

% --- TÍTULOS ---
\usepackage{titlesec}
\titleformat{\section}{\color{black}\large\bfseries}{\thesection}{0em}{\vspace{0.1cm}\titlerule[0.8pt]}
\titlespacing*{\section}{0pt}{10pt}{6pt} 
\titleformat{\subsubsection}[runin]{\color{black}\bfseries}{}{0em}{}
\titlespacing*{\subsubsection}{0pt}{0pt}{0.15cm}

% --- COLUNAS E GRÁFICOS ---
\usepackage{multicol}
\setlength{\columnsep}{1cm}
\setlength{\columnseprule}{0.5pt}
\def\columnseprulecolor{\color{CorLinha}}
\usepackage{adjustbox, tikz, pgfplots, circuitikz}
\usetikzlibrary{shadows, shapes.misc, positioning, calc}
\pgfplotsset{compat=1.18}

\begin{document}
\thispagestyle{empty}
\color{Cinzablack}
\vspace*{-1.8cm}

% --- CABEÇALHO PRINCIPAL AUTORAL (Apenas 1ª página) ---
\noindent
\begin{minipage}{\textwidth}
    \fontfamily{phv}\selectfont
    \noindent
    \begin{minipage}[t]{0.70\linewidth}
        {\Large \bfseries \MakeUppercase{Caderno de Atividades}}\\[0.1cm]
        {\large \textmd{\textcolor{CinzaEscuro}{Unidade: Decimais, Potenciação e Radiciação}}}
    \end{minipage}%
    \begin{minipage}[t]{0.30\linewidth}
        \raggedleft
        \small \textbf{Prof. Raphael Paulino}\\
        \textsf{\small \textcolor{indigo}{\texttt{profraphaelpaulino.com.br}}}
    \end{minipage}
    
    \vspace{0.15cm}
    \noindent\rule{\linewidth}{1.2pt}
    \vspace{-0.1cm}
    \footnotesize\textsf{\textbf{ÁREA:} MATEMÁTICA / ÁLGEBRA \hfill \textbf{NÍVEL:} EF II (7º ANO)}
    \vspace{0.1cm}
    \noindent\rule{\linewidth}{0.4pt}
\end{minipage}

\par\vspace{0.3cm}
\raggedcolumns
\begin{multicols*}{2}
\vspace{0.1cm}

\subsection*{Capítulo 1: Operações com Decimais}
\vspace{-0.2cm}\noindent\rule{\linewidth}{1.5pt} 
\par\vspace{\espacoPosTitulo}
\textbf{1.} Resolva a expressão numérica a seguir, respeitando a ordem de precedência das operações. Justifique o passo a passo da divisão entre números decimais, demonstrando como você igualou as casas decimais.
$$ \textbf{E} = (4,5 \times 1,2) - \left( 3,6 \div 0,4 \right) $$
\par\vspace{\espacoQuestao}

\noindent\textcolor{CorLinha}{\rule{\linewidth}{0.5pt}} 
\par\vspace{\espacoPreQuestao}
\textbf{2.} Julgue as afirmativas abaixo como Verdadeiras (V) ou Falsas (F). Para as falsas, reescreva a sentença corrigindo o erro matemático.
\begin{enumerate}[label=\Roman*.]
    \item Ao multiplicarmos $ 0,4 $ por $ 0,4 $, o resultado obtido é $ 1,6 $.
    \item Toda fração cujo denominador é uma potência de $ 10 $ pode ser escrita como um número decimal exato.
    \item A divisão de $ 7,5 $ por $ 0,1 $ produz o mesmo resultado que a multiplicação de $ 7,5 $ por $ 10 $.
\end{enumerate}
\par\vspace{\espacoQuestao}

\noindent\textcolor{CorLinha}{\rule{\linewidth}{0.5pt}} 
\par\vspace{\espacoPreQuestao}
\textbf{3.} Uma família decidiu trocar o piso retangular da sala de estar. O pedreiro enviou um esboço com as medidas do cômodo para que o proprietário pudesse calcular a quantidade de metros quadrados de cerâmica necessários.
\begin{center}
% ==========================================
% CONTROLE INDIVIDUAL DESTA IMAGEM:
% 1. CORTAR BORDAS: Mude os zeros do trim (Esquerda, Baixo, Direita, Topo). Ex: trim=1cm 0cm 1cm 0cm
% 2. TAMANHO: Para encolher só esta imagem, troque \larguraTikz por um valor (Ex: 0.5\linewidth)
% ==========================================
\begin{adjustbox}{trim=0cm 0cm 0cm 0cm, clip, max width=\larguraTikz}
\begin{tikzpicture}[x=1cm, y=1cm]
\definecolor{fundo}{HTML}{ECF0F1}
\definecolor{parede}{HTML}{2C3E50}
\definecolor{piso}{HTML}{E67E22}

% --- APLICANDO DROP SHADOW E ROUNDED CORNERS ---
% LAYER 1: Base/Sombra
\fill[fundo, drop shadow={opacity=0.2, shadow xshift=2pt, shadow yshift=-2pt}] (-0.5,-0.5) rectangle (5.5,3.5);

% LAYER 2: Piso da sala (laranja claro) e textura de grid leve
\fill[piso, rounded corners=2pt, opacity=0.7] (0,0) rectangle (5,3);
\draw[step=0.5cm, white, very thin, opacity=0.5] (0,0) grid (5,3);

% LAYER 3: Paredes
\draw[parede, line width=3pt, rounded corners=3pt] (0,0) rectangle (5,3);

% LAYER 4: Cotas (Com máscara para não cruzar linhas)
\draw[|<->|, >=latex, thick] (0, -0.3) -- (5, -0.3) node[midway, fill=fundo, inner sep=2pt, font=\bfseries] {5,4 m};
\draw[|<->|, >=latex, thick] (5.3, 0) -- (5.3, 3) node[midway, fill=fundo, inner sep=2pt, font=\bfseries, rotate=90] {3,5 m};
\end{tikzpicture}
\end{adjustbox}
\end{center}
Com base nas dimensões indicadas no esboço, calcule a área exata da sala em metros quadrados ($ \textbf{m}^{2} $). Arme a conta e demonstre o posicionamento da vírgula no resultado final.
\par\vspace{\espacoQuestao}

\noindent\textcolor{CorLinha}{\rule{\linewidth}{0.5pt}} 
\par\vspace{\espacoPreQuestao}
\textbf{4.} O estudante fictício Carlos tentou transformar a fração $ \mfrac{5}{8} $ em decimal e escreveu o seguinte raciocínio na lousa: 
\textit{"Como 8 não divide 5, eu inverto as posições e divido 8 por 5, o que dá 1,6. Logo, $ \mfrac{5}{8} = 1,6 $."}
\begin{boxresponda}
\textbf{Responda:} Diagnostique o erro conceitual cometido por Carlos e apresente o cálculo e o resultado corretos utilizando o algoritmo formal da divisão longa.
\end{boxresponda}
\par\vspace{\espacoQuestao}

\noindent\textcolor{CorLinha}{\rule{\linewidth}{0.5pt}} 
\par\vspace{\espacoPreQuestao}
\textbf{5.} Organize as ideias: elabore um fluxograma simples (em texto) ou mapa de regras explicando o que acontece com a posição da vírgula quando multiplicamos um número decimal por $ 100 $ e quando o dividimos por $ 10 $. Exemplifique sua explicação utilizando o número $ 42,8 $.
\par\vspace{\espacoQuestao}

\noindent\textcolor{CorLinha}{\rule{\linewidth}{0.5pt}} 
\par\vspace{\espacoPreQuestao}
\textbf{6.} Para realizar uma receita de bolo, Joana foi ao mercado e comprou $ \textbf{2,5 kg} $ de farinha de trigo, sendo que cada quilograma custa $ \textbf{R\$ 4,80} $. Ela pagou sua compra com uma nota de $ \textbf{R\$ 20,00} $. Determine o valor exato do troco que Joana deve receber.
\par\vspace{\espacoQuestao}

\noindent\textcolor{CorLinha}{\rule{\linewidth}{0.5pt}} 
\par\vspace{\espacoPreQuestao}
\textbf{7.} Na padaria do bairro, os clientes utilizam balanças digitais de autoatendimento. Um cliente colocou uma bandeja de pães de queijo sobre o equipamento e o visor digital imediatamente retornou os dados da pesagem e do preço.
\begin{center}
% ==========================================
% CONTROLE INDIVIDUAL DESTA IMAGEM:
% 1. CORTAR BORDAS: Mude os zeros do trim (Esquerda, Baixo, Direita, Topo). Ex: trim=1cm 0cm 1cm 0cm
% 2. TAMANHO: Para encolher só esta imagem, troque \larguraTikz por um valor (Ex: 0.5\linewidth)
% ==========================================
\begin{adjustbox}{trim=0cm 0cm 0cm 0cm, clip, max width=\larguraTikz}
\begin{tikzpicture}
\definecolor{balanca}{HTML}{BDC3C7}
\definecolor{prato}{HTML}{95A5A6}
\definecolor{visor}{HTML}{1A1A1A}
\definecolor{display}{HTML}{2ECC71}
\definecolor{texto}{HTML}{2C3E50}

% --- APLICANDO DROP SHADOW ---
% LAYER 1: Base da Balança
\fill[balanca, rounded corners=4pt, drop shadow={opacity=0.3, shadow xshift=2pt, shadow yshift=-2pt}] (0,0) rectangle (6, 2.5);

% LAYER 2: Prato de pesagem superior
\fill[prato, rounded corners=2pt] (0.5, 2.5) rectangle (5.5, 3.2);

% LAYER 3: Visores digitais negros
\fill[visor, rounded corners=2pt] (0.5, 0.5) rectangle (2.5, 1.8);
\fill[visor, rounded corners=2pt] (3.5, 0.5) rectangle (5.5, 1.8);

% LAYER 4: Textos nos visores e legendas
\node[text=display, font=\bfseries\Large] at (1.5, 1.1) {0,450};
\node[text=texto, font=\bfseries\footnotesize] at (1.5, 2.1) {Peso (kg)};

\node[text=display, font=\bfseries\Large] at (4.5, 1.1) {32,00};
\node[text=texto, font=\bfseries\footnotesize] at (4.5, 2.1) {Preço/kg (R\$)};
\end{tikzpicture}
\end{adjustbox}
\end{center}
Com base nas informações exibidas nos visores da balança, calcule o valor total a ser pago por essa bandeja de pães de queijo.
\par\vspace{\espacoQuestao}

\noindent\textcolor{CorLinha}{\rule{\linewidth}{0.5pt}} 
\par\vspace{\espacoPreQuestao}
\textbf{8.} Uma escola fará uma campanha de doação de álcool em gel. A direção comprou um galão contendo $ \textbf{15,5 litros} $ do produto e deseja fracioná-lo em pequenos frascos individuais. Sabendo que cada frasco tem capacidade exata para $ \textbf{0,25 litro} $, calcule quantos frascos poderão ser completamente enchidos.
\par\vspace{\espacoQuestao}

\noindent\textcolor{CorLinha}{\rule{\linewidth}{0.5pt}} 
\par\vspace{\espacoPreQuestao}
\textbf{9.} (Estilo UNICAMP - Adaptada) O gráfico de barras a seguir ilustra o índice pluviométrico diário (quantidade de chuva, em milímetros) registrado em uma cidade durante os quatro primeiros dias de uma semana.
\begin{center}
% ==========================================
% CONTROLE INDIVIDUAL DESTA IMAGEM:
% 1. CORTAR BORDAS: Mude os zeros do trim (Esquerda, Baixo, Direita, Topo). Ex: trim=1cm 0cm 1cm 0cm
% 2. TAMANHO: Para encolher só esta imagem, troque \larguraTikz por um valor (Ex: 0.5\linewidth)
% ==========================================
\begin{adjustbox}{trim=0cm 0cm 0cm 0cm, clip, max width=\larguraTikz}
\begin{tikzpicture}
\definecolor{fundoaxis}{HTML}{FAFAFA}
\definecolor{barra}{HTML}{3498DB}

% --- APLICANDO DROP SHADOW ---
\fill[fundoaxis, rounded corners=4pt, drop shadow={opacity=0.15, shadow xshift=2pt, shadow yshift=-2pt}] (-1.5,-1) rectangle (6,5);

\begin{axis}[
    ybar,
    width=7cm, height=5.5cm,
    axis lines=middle,
    ymin=0, ymax=18,
    xmin=0, xmax=5,
    xtick={1,2,3,4},
    xticklabels={Seg, Ter, Qua, Qui},
    ytick={0, 5, 10, 15},
    ylabel={Chuva (mm)},
    nodes near coords,
    nodes near coords style={font=\bfseries, text=black, yshift=4pt},
    bar width=20pt,
    enlarge x limits=0.15
]
% LAYER: Barras
\addplot[fill=barra, rounded corners=1pt] coordinates {
    (1, 12.4)
    (2, 8.5)
    (3, 14.8)
    (4, 5.3)
};
\end{axis}
\end{tikzpicture}
\end{adjustbox}
\end{center}
Determine a diferença exata, em milímetros, entre o dia em que mais choveu e o dia em que menos choveu nesse período. Além de assinalar a alternativa correta, \textbf{apresente o cálculo refutando matematicamente o erro contido na alternativa (a)}.
\begin{enumerate}[label=\alph*)]
    \item $ 14,8 \text{ mm} $
    \item $ 9,5 \text{ mm} $
    \item $ 7,1 \text{ mm} $
    \item $ 3,9 \text{ mm} $
\end{enumerate}
\par\vspace{\espacoQuestao}

\noindent\textcolor{CorLinha}{\rule{\linewidth}{0.5pt}} 
\par\vspace{\espacoPreQuestao}
\textbf{10.} (ENEM - Adaptada) Um motorista de aplicativo abasteceu seu veículo com gasolina. A bomba registrou que foram colocados exatos $ \textbf{40,5 litros} $ no tanque. Se o preço do litro de gasolina no posto era de $ \textbf{R\$ 5,40} $, assinale a alternativa que apresenta o valor total pago por esse abastecimento.
\begin{enumerate}[label=\alph*)]
    \item $ \text{R\$} \ 218,70 $
    \item $ \text{R\$} \ 228,50 $
    \item $ \text{R\$} \ 216,40 $
    \item $ \text{R\$} \ 214,70 $
\end{enumerate}
\par\vspace{\espacoQuestao}

\subsection*{Capítulo 2: Potenciação de Racionais}
\vspace{-0.2cm}\noindent\rule{\linewidth}{1.5pt} 
\par\vspace{\espacoPosTitulo}
\textbf{11.} A potenciação no conjunto dos números racionais obedece a regras de sinais dependendo do expoente. Calcule o valor exato da expressão algébrica a seguir, demonstrando os resultados parciais antes da soma.
$$ \textbf{P} = \left( -\mfrac{1}{2} \right)^{3} + \left( \mfrac{3}{4} \right)^{2} $$
\par\vspace{\espacoQuestao}

\noindent\textcolor{CorLinha}{\rule{\linewidth}{0.5pt}} 
\par\vspace{\espacoPreQuestao}
\textbf{12.} Julgue as afirmativas matemáticas a seguir como Verdadeiras (V) ou Falsas (F). Justifique as sentenças que julgar incorretas.
\begin{enumerate}[label=\Roman*.]
    \item Toda potência de base negativa, elevada a um expoente ímpar, resulta em um número racional negativo.
    \item O cálculo de $ \left( \mfrac{2}{5} \right)^{-2} $ equivale à inversão da base, resultando na fração $ \mfrac{25}{4} $.
    \item A potência $ \left( 0,3 \right)^{2} $ tem como resultado exato o número decimal $ 0,9 $.
\end{enumerate}
\par\vspace{\espacoQuestao}

\noindent\textcolor{CorLinha}{\rule{\linewidth}{0.5pt}} 
\par\vspace{\espacoPreQuestao}
\textbf{13.} Durante uma aula prática de geometria e arte, um aluno recorta peças quadradas de papelão para montar um mosaico. Ele recebe uma placa matriz e precisa calcular a área exata da superfície.
\par\vspace{0.1cm}
\begin{center}
% ==========================================
% CONTROLE INDIVIDUAL DESTA IMAGEM:
% 1. CORTAR BORDAS: Mude os zeros do trim (Esquerda, Baixo, Direita, Topo). Ex: trim=1cm 0cm 1cm 0cm
% 2. TAMANHO: Para encolher só esta imagem, troque \larguraTikz por um valor (Ex: 0.5\linewidth)
% ==========================================
\begin{adjustbox}{trim=1.2cm 1.2cm 1.2cm 1.2cm, clip, max width=0.6\linewidth}
\begin{tikzpicture}
\definecolor{papelao}{HTML}{D4AC0D}
\definecolor{corte}{HTML}{9A7D0A}
\definecolor{fundo}{HTML}{FDFEFE}

% --- APLICANDO DROP SHADOW ---
% LAYER 1: Bancada
\fill[fundo, rounded corners=2pt, drop shadow={opacity=0.2, shadow xshift=2pt, shadow yshift=-2pt}] (-1, -1) rectangle (6, 5);

% LAYER 2: Placa Quadrada
\fill[papelao, rounded corners=1pt, drop shadow={opacity=0.4}] (1, 1) rectangle (4, 4);
\draw[corte, thick, dashed] (1,1) rectangle (4,4);

% LAYER 3: Indicação de medidas (Fração)
\draw[|<->|, >=latex, thick] (1, 0.5) -- (4, 0.5) node[midway, fill=fundo, inner sep=2pt, font=\bfseries] {$ \mfrac{5}{2} $ cm};
\draw[|<->|, >=latex, thick] (4.7, 1) -- (4.7, 4) node[midway, fill=fundo, inner sep=2pt, font=\bfseries] {$ \mfrac{5}{2} $ cm};
\end{tikzpicture}
\end{adjustbox}
\end{center}
Exiba a montagem da equação da área utilizando a notação de potenciação (base e expoente) e, em seguida, resolva-a, entregando o resultado final em formato de fração irredutível.
\par\vspace{\espacoQuestao}

\noindent\textcolor{CorLinha}{\rule{\linewidth}{0.5pt}} 
\par\vspace{\espacoPreQuestao}
\textbf{14.} O aluno Marcelo estava resolvendo uma tarefa sobre propriedades das potências e simplificou a expressão fracionária abaixo da seguinte maneira:
$$
\begin{aligned}
& \textbf{M} = \mfrac{ \left( \mfrac{3}{7} \right)^{4} \cdot \left( \mfrac{3}{7} \right)^{5} }{ \left( \mfrac{3}{7} \right)^{2} }
\end{aligned}
$$
\par\vspace{0.2em}
Marcelo fez:
\par\vspace{0.2em}
$$
\begin{aligned}
& \mfrac{3}{7}^{(4 \cdot 5)} \div \mfrac{3}{7}^{2} = \mfrac{3}{7}^{20 - 2} = \left( \mfrac{3}{7} \right)^{18}
\end{aligned}
$$

\begin{boxresponda}
\textbf{Responda:} Identifique e explique textualmente a propriedade matemática que Marcelo confundiu no numerador. Em seguida, refaça o cálculo simplificando a expressão corretamente até reduzi-la a uma única potência.
\end{boxresponda}
\par\vspace{\espacoQuestao}
\noindent\textcolor{CorLinha}{\rule{\linewidth}{0.5pt}} 
\par\vspace{\espacoPreQuestao}
\textbf{15.} Aplica-se o conceito de expoente negativo invertendo-se a base racional. Resolva a expressão algébrica abaixo passo a passo, exibindo a transformação dos expoentes negativos em positivos e determinando o resultado numérico final em forma de fração irredutível.
$$ \textbf{K} = \left( \mfrac{3}{2} \right)^{-2} + \left( -\mfrac{1}{3} \right)^{2} $$
\par\vspace{\espacoQuestao}

\noindent\textcolor{CorLinha}{\rule{\linewidth}{0.5pt}} 
\par\vspace{\espacoPreQuestao}
\textbf{16.} Um projeto de artesanato utiliza chapas quadradas de acrílico para montar luminárias. O fornecedor enviou uma chapa piloto com as especificações dimensionais exibidas na planta baixa.
\par\vspace{0.3cm}
\begin{center}
% ==========================================
% CONTROLE INDIVIDUAL DESTA IMAGEM:
% 1. CORTAR BORDAS: Mude os zeros do trim (Esquerda, Baixo, Direita, Topo). Ex: trim=1cm 0cm 1cm 0cm
% 2. TAMANHO: Para encolher só esta imagem, troque \larguraTikz por um valor (Ex: 0.5\linewidth)
% ==========================================
\begin{adjustbox}{trim=0cm 0cm 0cm 0cm, clip, max width=\larguraTikz}
\begin{tikzpicture}
\definecolor{acrilico}{HTML}{3498DB}
\definecolor{fundo}{HTML}{ECF0F1}
\definecolor{cota}{HTML}{2C3E50}

% --- APLICANDO DROP SHADOW ---
% LAYER 1: Bancada
%\fill[fundo, drop shadow={opacity=0.2, shadow xshift=2pt, shadow yshift=-2pt}, rounded corners=2pt] (-1, -1) rectangle (5, 5);

% LAYER 2: Chapa de Acrílico (Quadrada com opacidade para vidro)
\fill[acrilico, opacity=0.6, drop shadow={opacity=0.4, shadow xshift=3pt, shadow yshift=-3pt}, rounded corners=1pt] (0.5, 0.5) rectangle (3.5, 3.5);
\draw[acrilico!80!black, thick, rounded corners=1pt] (0.5, 0.5) rectangle (3.5, 3.5);

% LAYER 3: Cotas (Com máscara de leitura branca)
\draw[|<->|, >=latex, thick, cota] (0.5, 0.1) -- (3.5, 0.1) node[midway, fill=fundo, inner sep=2pt, font=\bfseries] {1,4 m};
\draw[|<->|, >=latex, thick, cota] (3.9, 0.5) -- (3.9, 3.5) node[midway, fill=fundo, inner sep=2pt, font=\bfseries, rotate=90] {1,4 m};
\end{tikzpicture}
\end{adjustbox}
\end{center}
Modele a equação da área da superfície dessa chapa de acrílico utilizando o conceito de potenciação e, em seguida, calcule o resultado exato em metros quadrados.
\par\vspace{\espacoQuestao}

\noindent\textcolor{CorLinha}{\rule{\linewidth}{0.5pt}} 
\par\vspace{\espacoPreQuestao}
\textbf{17.} O estudante fictício João estava resolvendo exercícios sobre base negativa e escreveu a seguinte afirmação na lousa: 
\textit{"As expressões $ -0,5^{2} $ e $ (-0,5)^{2} $ produzem exatamente o mesmo resultado numérico, pois o expoente é par em ambas, o que sempre garante um resultado positivo, ou seja, $ +0,25 $."}
\begin{boxresponda}
\textbf{Responda:} João cometeu um erro grave de sintaxe matemática e de interpretação da regra de sinais. Diagnostique e explique o erro textualmente. Em seguida, prove matematicamente calculando o valor exato de cada uma das duas expressões.
\end{boxresponda}
\par\vspace{\espacoQuestao}

\noindent\textcolor{CorLinha}{\rule{\linewidth}{0.5pt}} 
\par\vspace{\espacoPreQuestao}
\textbf{18.} (Estilo UNICAMP) Um biólogo monitora o crescimento de uma cultura de bactérias em laboratório. A população inicial era de apenas $ \textbf{1 bactéria} $. O pesquisador observou que, devido às condições ideais da estufa, o número de bactérias triplica exatamente a cada hora. Modele a expressão em forma de potência que representa a quantidade de bactérias após exatas $ \textbf{5 horas} $ e calcule o número total de micro-organismos nesse instante.
\par\vspace{\espacoQuestao}

\noindent\textcolor{CorLinha}{\rule{\linewidth}{0.5pt}} 
\par\vspace{\espacoPreQuestao}
\textbf{19.} Um processo mecânico de filtragem de água ocorre em múltiplos estágios. A cada passagem por um filtro cilíndrico, a quantidade de partículas sólidas na água é reduzida estritamente à metade. Sabendo que a água passa por $ \textbf{4 filtros} $ consecutivos, assinale a alternativa que traz a fração exata de partículas que restará na água. Além de achar a certa, \textbf{refute matematicamente a alternativa (a), explicando que erro comum de potenciação leva o aluno a escolhê-la.}
\begin{enumerate}[label=\alph*)]
    \item $ \mfrac{1}{8} $
    \item $ \mfrac{1}{16} $
    \item $ \mfrac{1}{32} $
    \item $ \mfrac{1}{4} $
\end{enumerate}
\par\vspace{\espacoQuestao}

\noindent\textcolor{CorLinha}{\rule{\linewidth}{0.5pt}} 
\par\vspace{\espacoPreQuestao}
\textbf{20.} Um arquiteto de redes estruturou o cabeamento de um pequeno prédio comercial utilizando um padrão fixo de ramificação exponencial, partindo do servidor em direção aos computadores.
\begin{center}
% ==========================================
% CONTROLE INDIVIDUAL DESTA IMAGEM:
% 1. CORTAR BORDAS: Mude os zeros do trim (Esquerda, Baixo, Direita, Topo). Ex: trim=1cm 0cm 1cm 0cm
% 2. TAMANHO: Para encolher só esta imagem, troque \larguraTikz por um valor (Ex: 0.5\linewidth)
% ==========================================
\begin{adjustbox}{trim=0cm 0cm 0cm 0cm, clip, max width=\larguraTikz}
\begin{tikzpicture}
\definecolor{servidor}{HTML}{8E44AD}
\definecolor{roteador}{HTML}{3498DB}
\definecolor{fundo}{HTML}{FDFEFE}
\definecolor{linha}{HTML}{7F8C8D}

% --- APLICANDO DROP SHADOW ---
% LAYER 1: Background
\fill[fundo, rounded corners=3pt, drop shadow={opacity=0.2}] (-1, -0.5) rectangle (9, 4.5);

% LAYER 2: Conexões (Linhas)
\foreach \x in {0.5, 4, 7.5} {
    \draw[linha, line width=1.5pt] (4, 3.5) -- (\x, 1.5);
}

% LAYER 3: Nós da Rede
% Servidor (Nível 0)
\fill[servidor, rounded corners=2pt, drop shadow={opacity=0.3, shadow xshift=2pt, shadow yshift=-2pt}] (3.2, 3.3) rectangle (4.8, 3.9);
\node[text=white, font=\bfseries] at (4, 3.6) {Servidor};

% Roteadores (Nível 1)
\foreach \x in {0.5, 4, 7.5} {
    \fill[roteador, rounded corners=2pt, drop shadow={opacity=0.3, shadow xshift=2pt, shadow yshift=-2pt}] (\x-1, 1.2) rectangle (\x+1, 1.8);
    \node[text=white, font=\bfseries\small] at (\x, 1.5) {Roteador};
}

% Labels de ramificação
\node[fill=fundo, text=linha, inner sep=2pt, font=\bfseries\small] at (4, 2.5) {3 Conexões};
\node[fill=fundo, text=linha, inner sep=2pt, font=\bfseries] at (4, 0.5) {Cada roteador liga a 3 PCs};
\end{tikzpicture}
\end{adjustbox}
\end{center}
Com base na topologia da rede ilustrada, escreva a potência matemática que representa o total de computadores (PCs) no final do processo e calcule esse número exato de estações de trabalho.
\par\vspace{\espacoQuestao}

\noindent\textcolor{CorLinha}{\rule{\linewidth}{0.5pt}} 
\par\vspace{\espacoPreQuestao}
\textbf{21.} (ENEM - Adaptada) O origami é a arte tradicional japonesa de dobrar o papel. Um estudante pegou uma folha de papel cuja espessura original era de $ \textbf{0,1 mm} $ e começou a dobrá-la exatamente ao meio, sucessivas vezes. Calcule a espessura total, em milímetros, que a folha dobrada atingirá após ser dobrada ao meio $ \textbf{6 vezes} $ consecutivas.
\par\vspace{\espacoQuestao}

\noindent\textcolor{CorLinha}{\rule{\linewidth}{0.5pt}} 
\par\vspace{\espacoPreQuestao}
\textbf{22.} Em um torneio de xadrez escolar eliminatório, metade dos competidores é eliminada a cada rodada. O torneio iniciou com exatos $ \textbf{128 alunos} $. O estudante fictício Marcos tentou calcular quantos jogadores restariam após a quinta rodada fazendo a seguinte operação: $ 128 \div 5 = 25,6 $ alunos.
\begin{boxresponda}
\textbf{Responda:} Diagnostique o absurdo lógico e matemático na modelagem de Marcos. Em seguida, utilize a multiplicação por potências de fração para demonstrar o cálculo correto e encontrar o número real de alunos que restarão.
\end{boxresponda}
\par\vspace{\espacoQuestao}

\subsection*{Capítulo 3: Raiz Quadrada Exata}
\vspace{-0.2cm}\noindent\rule{\linewidth}{1.5pt} 
\par\vspace{\espacoPosTitulo}
\textbf{23.} Para extrair a raiz quadrada exata de uma fração, devemos operar separadamente o numerador e o denominador. Com base nessa premissa fundamental, calcule o valor da expressão numérica a seguir, entregando o resultado final de forma simplificada.
$$ \textbf{R} = \sqrt{\mfrac{81}{144}} + \mfrac{1}{4} $$
\par\vspace{\espacoQuestao}

\noindent\textcolor{CorLinha}{\rule{\linewidth}{0.5pt}} 
\par\vspace{\espacoPreQuestao}
\textbf{24.} A direção de uma escola contratou uma empresa de engenharia para demarcar um novo pátio cimentado no formato de um quadrado perfeito.
\begin{center}
% ==========================================
% CONTROLE INDIVIDUAL DESTA IMAGEM:
% 1. CORTAR BORDAS: Mude os zeros do trim (Esquerda, Baixo, Direita, Topo). Ex: trim=1cm 0cm 1cm 0cm
% 2. TAMANHO: Para encolher só esta imagem, troque \larguraTikz por um valor (Ex: 0.5\linewidth)
% ==========================================
\begin{adjustbox}{trim=0cm 0cm 0cm 0cm, clip, max width=\larguraTikz}
\begin{tikzpicture}
\definecolor{patio}{HTML}{95A5A6}
\definecolor{grama}{HTML}{27AE60}
\definecolor{cota}{HTML}{2C3E50}
\definecolor{texto}{HTML}{FFFFFF}

% --- APLICANDO DROP SHADOW ---
% LAYER 1: Gramado de fundo
\fill[grama, rounded corners=3pt, drop shadow={opacity=0.2, shadow xshift=2pt, shadow yshift=-2pt}] (-1, -1) rectangle (6, 6);
\draw[step=1cm, black, very thin, opacity=0.1] (-1,-1) grid (6,6);

% LAYER 2: Pátio Quadrado
\fill[patio, rounded corners=1pt, drop shadow={opacity=0.5, shadow xshift=2pt, shadow yshift=-2pt}] (1, 1) rectangle (4, 4);
\draw[white, thick, dashed] (1,1) rectangle (4,4);

% LAYER 3: Identificação de Área (Com máscara)
\node[fill=grama, text=texto, inner sep=3pt, font=\bfseries] at (2.5, 2.5) {Área = 20,25 m²};

% LAYER 4: Identificação de Lado
\draw[|<->|, >=latex, thick, white] (1, 0.5) -- (4, 0.5) node[midway, fill=grama, inner sep=2pt, font=\bfseries] {Lado \textbf{x}};
\end{tikzpicture}
\end{adjustbox}
\end{center}
Transforme a área informada na planta para a forma de fração decimal e extraia a raiz quadrada para encontrar a medida exata do lado $ \textbf{x} $ desse pátio, em metros.
\par\vspace{\espacoQuestao}

\noindent\textcolor{CorLinha}{\rule{\linewidth}{0.5pt}} 
\par\vspace{\espacoPreQuestao}
\textbf{25.} Julgue as afirmativas matemáticas a respeito da radiciação, utilizando Verdadeiro (V) ou Falso (F). Apresente a correção matemática para as sentenças falsas.
\begin{enumerate}[label=\Roman*.]
    \item A operação de raiz quadrada atua como o processo inverso da potenciação de expoente $ 2 $.
    \item No conjunto dos números racionais, a raiz quadrada exata de $ -25 $ tem como resultado o número $ -5 $.
    \item A raiz quadrada exata do número decimal $ 0,64 $ corresponde estritamente à fração $ \mfrac{4}{5} $.
\end{enumerate}
\par\vspace{\espacoQuestao}

\begin{boxdica}
\textbf{Dica Metodológica:} Para extrair a raiz de decimais pequenos sem errar a posição da vírgula, converta-os para fração com denominador $ 100 $, $ 10.000 $, etc., tire a raiz de cima e de baixo separadamente, e depois volte para decimal!
\end{boxdica}

\noindent\textcolor{CorLinha}{\rule{\linewidth}{0.5pt}} 
\par\vspace{\espacoPreQuestao}
\textbf{26.} Aplicando o raciocínio indicado na dica anterior, demonstre o passo a passo operatório para calcular o valor exato da expressão $ \sqrt{0,0144} $. Exiba o resultado final obrigatoriamente no formato decimal.
\par\vspace{\espacoQuestao}

\noindent\textcolor{CorLinha}{\rule{\linewidth}{0.5pt}} 
\par\vspace{\espacoPreQuestao}
\textbf{27.} O estudante fictício Bruno tentou calcular o tamanho da hipotenusa de um triângulo retângulo e chegou na seguinte equação: $ \textbf{h} = \sqrt{16 + 9} $. Ao resolver, ele anotou: 
\textit{"Eu sei que a raiz de 16 é 4 e a raiz de 9 é 3. Como é uma soma, o resultado de h é $ 4 + 3 = 7 $."}
\begin{boxresponda}
\textbf{Responda:} Bruno cometeu um dos erros mais clássicos da radiciação. Explique teoricamente o porquê da regra usada por ele estar errada, calcule o resultado exato correto de $ \textbf{h} $ e prove numericamente a diferença.
\end{boxresponda}
\par\vspace{\espacoQuestao}

\noindent\textcolor{CorLinha}{\rule{\linewidth}{0.5pt}} 
\par\vspace{\espacoPreQuestao}
\textbf{28.} Uma vidraçaria corta tampos de vidro temperado estritamente no formato de quadrados perfeitos. O controle de qualidade inspeciona duas peças produzidas em sequência.
\begin{center}
% ==========================================
% CONTROLE INDIVIDUAL DESTA IMAGEM:
% 1. CORTAR BORDAS: Mude os zeros do trim (Esquerda, Baixo, Direita, Topo). Ex: trim=1cm 0cm 1cm 0cm
% 2. TAMANHO: Para encolher só esta imagem, troque \larguraTikz por um valor (Ex: 0.5\linewidth)
% ==========================================
\begin{adjustbox}{trim=0cm 1cm 0cm 0.7cm, clip, max width=\larguraTikz}
\begin{tikzpicture}
\definecolor{bancada}{HTML}{E5E8E8}
\definecolor{vidroA}{HTML}{3498DB}
\definecolor{vidroB}{HTML}{2980B9}
\definecolor{etiqueta}{HTML}{FFFFFF}
\definecolor{texto}{HTML}{2C3E50}

% --- APLICANDO DROP SHADOW ---
% LAYER 1: Bancada
\fill[bancada, rounded corners=3pt, drop shadow={opacity=0.2, shadow xshift=2pt, shadow yshift=-2pt}] (-1, -1) rectangle (8, 5);

% LAYER 2: Tampos de Vidro
\fill[vidroA, opacity=0.7, rounded corners=1pt, drop shadow={opacity=0.4, shadow xshift=3pt, shadow yshift=-3pt}] (0.5, 0.5) rectangle (3.5, 3.5);
\fill[vidroB, opacity=0.7, rounded corners=1pt, drop shadow={opacity=0.4, shadow xshift=3pt, shadow yshift=-3pt}] (4.5, 0.5) rectangle (6.5, 2.5);

% LAYER 3: Brilho/Reflexo
\draw[white, thick, opacity=0.5] (0.8, 3.2) -- (1.5, 3.2);
\draw[white, thick, opacity=0.5] (4.8, 2.2) -- (5.5, 2.2);

% LAYER 4: Etiquetas (Máscaras de leitura)
\node[fill=etiqueta, text=texto, inner sep=3pt, font=\bfseries] at (2, 2) {Área = 0,81 m²};
\node[fill=etiqueta, text=texto, inner sep=3pt, font=\bfseries, text width=2cm, align=center] at (5.5, 1.5) {Área = 1,21 m²};
\node[text=texto, font=\bfseries\small] at (2, 4) {Tampo A};
\node[text=texto, font=\bfseries\small] at (5.5, 3) {Tampo B};
\end{tikzpicture}
\end{adjustbox}
\end{center}
Determine o comprimento exato do lado de cada tampo de vidro, calculando as respectivas raízes. Em seguida, determine a diferença, em metros, entre as medidas.
\par\vspace{\espacoQuestao}
\noindent\textcolor{CorLinha}{\rule{\linewidth}{0.5pt}} 
\par\vspace{\espacoPreQuestao}
\textbf{29.} (Estilo VUNESP) O tabuleiro do clássico jogo de xadrez é um grande quadrado perfeito formado por exatamente $ \textbf{64 casas} $ menores, também quadradas e idênticas, intercaladas em cores claras e escuras. Sabendo que a área total da superfície desse tabuleiro mede exatos $ \textbf{1.024 cm}^{2} $, determine a medida do lado de cada uma das pequenas casas quadradas do jogo.
\par\vspace{\espacoQuestao}

\noindent\textcolor{CorLinha}{\rule{\linewidth}{0.5pt}} 
\par\vspace{\espacoPreQuestao}
\textbf{30.} Um professor construiu um varal matemático em sala de aula para representar a reta numérica e marcou posições específicas, exigindo dos alunos a alocação correta de algumas raízes exatas.
\begin{center}
% ==========================================
% CONTROLE INDIVIDUAL DESTA IMAGEM:
% 1. CORTAR BORDAS: Mude os zeros do trim (Esquerda, Baixo, Direita, Topo). Ex: trim=1cm 0cm 1cm 0cm
% 2. TAMANHO: Para encolher só esta imagem, troque \larguraTikz por um valor (Ex: 0.5\linewidth)
% ==========================================
\begin{adjustbox}{trim=0cm 0cm 0cm 1.2cm, clip, max width=\larguraTikz}
\begin{tikzpicture}
\definecolor{reta}{HTML}{34495E}
\definecolor{ponto}{HTML}{E74C3C}
\definecolor{fundo}{HTML}{FDFEFE}

% --- APLICANDO DROP SHADOW ---
% LAYER 1: Fundo da Régua
\fill[fundo, rounded corners=2pt, drop shadow={opacity=0.15, shadow xshift=2pt, shadow yshift=-2pt}] (-1, -1) rectangle (10, 2);

% LAYER 2: Eixo Principal
\draw[->, >=latex, line width=2pt, reta] (0,0) -- (9,0);

% LAYER 3: Marcações Inteiras e Meios
\foreach \x/\label in {1/0, 3/1, 5/2, 7/3} {
    \draw[line width=1.5pt, reta] (\x, -0.2) -- (\x, 0.2);
    \node[font=\bfseries, text=reta] at (\x, -0.6) {\label};
}
\foreach \x in {2, 4, 6, 8} {
    \draw[line width=1pt, reta] (\x, -0.1) -- (\x, 0.1);
}

% LAYER 4: Pontos A, B, C (Com máscaras)
\fill[ponto, drop shadow={opacity=0.4}] (4, 0) circle (4pt) node[above=0.2cm, fill=fundo, inner sep=1pt, font=\bfseries] {Ponto A};
\fill[ponto, drop shadow={opacity=0.4}] (5, 0) circle (4pt) node[above=0.2cm, fill=fundo, inner sep=1pt, font=\bfseries] {Ponto B};
\fill[ponto, drop shadow={opacity=0.4}] (6, 0) circle (4pt) node[above=0.2cm, fill=fundo, inner sep=1pt, font=\bfseries] {Ponto C};
\end{tikzpicture}
\end{adjustbox}
\end{center}
Dada a expressão numérica $ \textbf{x} = \sqrt{2,25} $, calcule o valor exato de $ \textbf{x} $ e indique qual dos pontos demarcados na reta (A, B ou C) corresponde rigorosamente a esse resultado.
\par\vspace{\espacoQuestao}

\noindent\textcolor{CorLinha}{\rule{\linewidth}{0.5pt}} 
\par\vspace{\espacoPreQuestao}
\textbf{31.} (ENEM - Adaptada) O proprietário de uma fazenda deseja cercar completamente um curral que possui o formato geométrico de um quadrado perfeito. O engenheiro agrônomo informa que a área total destinada a esse curral mede exatamente $ \textbf{256 m}^{2} $. Sabendo que será instalada uma cerca utilizando $ \textbf{4 fios} $ de arame, assinale a alternativa que indica a quantidade total de arame que deverá ser comprada. \textbf{Além de marcar a resposta correta, diagnostique e refute matematicamente o distrator presente na alternativa (b).}
\begin{enumerate}[label=\alph*)]
    \item $ 16 \text{ m} $
    \item $ 64 \text{ m} $
    \item $ 128 \text{ m} $
    \item $ 256 \text{ m} $
\end{enumerate}
\par\vspace{\espacoQuestao}

\noindent\textcolor{CorLinha}{\rule{\linewidth}{0.5pt}} 
\par\vspace{\espacoPreQuestao}
\textbf{32.} O estudante fictício Ricardo tentou calcular raízes de números decimais baseando-se apenas nos algarismos significativos. Ele escreveu na prova: 
\textit{"Como $ \sqrt{144} = 12 $, então a raiz de $ \textbf{14,4} $ é obviamente $ \textbf{1,2} $."}
\begin{boxresponda}
\textbf{Responda:} Ricardo caiu em uma famosa "pegadinha" da radiciação decimal. Explique matematicamente por que $ \sqrt{14,4} $ NÃO é igual a $ 1,2 $, demonstrando qual seria a fração geratriz correta se ele quisesse que a resposta fosse $ 1,2 $.
\end{boxresponda}
\par\vspace{\espacoQuestao}

\noindent\textcolor{CorLinha}{\rule{\linewidth}{0.5pt}} 
\par\vspace{\espacoPreQuestao}
\textbf{33.} A malha quadriculada apresenta a planta baixa da sala de testes de um laboratório de robótica. A área total reservada para o circuito dos robôs é composta pela união de diversos quadrados idênticos pintados em azul.
\begin{center}
% ==========================================
% CONTROLE INDIVIDUAL DESTA IMAGEM:
% 1. CORTAR BORDAS: Mude os zeros do trim (Esquerda, Baixo, Direita, Topo). Ex: trim=1cm 0cm 1cm 0cm
% 2. TAMANHO: Para encolher só esta imagem, troque \larguraTikz por um valor (Ex: 0.5\linewidth)
% ==========================================
\begin{adjustbox}{trim=0cm 0cm 0cm 0cm, clip, max width=\larguraTikz}
\begin{tikzpicture}
\definecolor{malha}{HTML}{BDC3C7}
\definecolor{circuito}{HTML}{8E44AD}
\definecolor{fundo}{HTML}{FAFAFA}
\definecolor{texto}{HTML}{2C3E50}

% --- APLICANDO DROP SHADOW ---
% LAYER 1: Base e Malha
\fill[fundo, drop shadow={opacity=0.2, shadow xshift=2pt, shadow yshift=-2pt}, rounded corners=2pt] (-0.5, -0.5) rectangle (5.5, 4.5);
\draw[step=1cm, malha, very thin] (0,0) grid (5,4);

% LAYER 2: Quadrados do circuito
\foreach \x/\y in {1/1, 2/1, 3/1, 2/2, 2/3} {
    \fill[circuito, rounded corners=1pt, drop shadow={opacity=0.3, shadow xshift=1pt}] (\x+0.1, \y+0.1) rectangle (\x+0.9, \y+0.9);
}

% LAYER 3: Detalhe de área
\draw[->, >=latex, thick, texto] (4.5, 3.5) -- (2.8, 3.5);
\node[fill=fundo, text=texto, inner sep=2pt, font=\bfseries\small] at (4.5, 3.8) {Área Unitária};
\node[fill=fundo, text=texto, inner sep=2pt, font=\bfseries\small] at (4.5, 3.2) {= $ 0,25 \text{ m}^{2} $};
\end{tikzpicture}
\end{adjustbox}
\end{center}
Baseando-se na área de cada um dos quadrados azuis isoladamente, calcule a medida do lado de um único quadrado. Em seguida, determine a medida total do perímetro (contorno externo) dessa pista de robótica.
\par\vspace{\espacoQuestao}

\noindent\textcolor{CorLinha}{\rule{\linewidth}{0.5pt}} 
\par\vspace{\espacoPreQuestao}
\textbf{34.} O cálculo da raiz quadrada de números grandes pode ser facilitado aplicando-se a fatoração (decomposição em fatores primos). Utilizando esse método rigorosamente (agrupamento de pares iguais), extraia passo a passo a raiz quadrada de $ \textbf{1.296} $.
\par\vspace{\espacoQuestao}
\noindent\textcolor{CorLinha}{\rule{\linewidth}{0.5pt}} 
\par\vspace{\espacoPreQuestao}
\textbf{35.} (Estilo UNICAMP - Adaptada) Uma indústria produz painéis solares no formato de grandes quadrados perfeitos. Por razões de segurança aerodinâmica, a engenharia estabelece que a área útil de captação de um modelo específico de painel deve ser de exatamente $ \textbf{3,24 m}^{2} $. Determine o comprimento do lado desse painel em metros.
\par\vspace{\espacoQuestao}

\noindent\textcolor{CorLinha}{\rule{\linewidth}{0.5pt}} 
\par\vspace{\espacoPreQuestao}
\textbf{36.} Calcule o valor exato da expressão algébrica abaixo, que exige o domínio da extração de raízes fracionárias seguida de uma multiplicação. Simplifique o resultado obtido até transformá-lo em uma fração irredutível.
$$ \textbf{M} = \sqrt{\mfrac{25}{36}} \cdot \sqrt{\mfrac{49}{100}} $$
\par\vspace{\espacoQuestao}

\subsection*{Capítulo 4: Razão e Proporção - Introdução}
\vspace{-0.2cm}\noindent\rule{\linewidth}{1.5pt} 
\par\vspace{\espacoPosTitulo}
\textbf{37.} No estacionamento de um shopping, um sensor fotográfico cataloga os veículos de acordo com a sua tonalidade de cor para gerar estatísticas no painel de controle.
\begin{center}
% ==========================================
% CONTROLE INDIVIDUAL DESTA IMAGEM:
% 1. CORTAR BORDAS: Mude os zeros do trim (Esquerda, Baixo, Direita, Topo). Ex: trim=1cm 0cm 1cm 0cm
% 2. TAMANHO: Para encolher só esta imagem, troque \larguraTikz por um valor (Ex: 0.5\linewidth)
% ==========================================
\begin{adjustbox}{trim=0cm 0cm 0cm 0cm, clip, max width=\larguraTikz}
\begin{tikzpicture}
\definecolor{asfalto}{HTML}{7F8C8D}
\definecolor{carrovermelho}{HTML}{E74C3C}
\definecolor{carroazul}{HTML}{3498DB}
\definecolor{carroprata}{HTML}{ECF0F1}
\definecolor{fundo}{HTML}{FDFEFE}
\definecolor{texto}{HTML}{2C3E50}

% --- APLICANDO DROP SHADOW ---
% LAYER 1: Estacionamento
\fill[asfalto, rounded corners=4pt, drop shadow={opacity=0.3, shadow xshift=2pt, shadow yshift=-2pt}] (-1, 0) rectangle (7, 4);

% Marcações de Vagas
\foreach \x in {0, 1.5, 3, 4.5, 6} {
    \draw[white, line width=2pt] (\x, 0.5) -- (\x, 3.5);
}

% LAYER 2: Carros (Retângulos com cantos arredondados simulando teto e capô)
% Vaga 1 (Prata)
\fill[carroprata, rounded corners=2pt, drop shadow={opacity=0.5}] (0.2, 1) rectangle (1.3, 3);
% Vaga 2 (Vermelho)
\fill[carrovermelho, rounded corners=2pt, drop shadow={opacity=0.5}] (1.7, 1) rectangle (2.8, 3);
% Vaga 3 (Azul)
\fill[carroazul, rounded corners=2pt, drop shadow={opacity=0.5}] (3.2, 1) rectangle (4.3, 3);
% Vaga 4 (Vermelho)
\fill[carrovermelho, rounded corners=2pt, drop shadow={opacity=0.5}] (4.7, 1) rectangle (5.8, 3);

% LAYER 3: Legenda do painel
\fill[fundo, rounded corners=2pt, drop shadow={opacity=0.4}] (1, -1.2) rectangle (5.7, -0.2);
\node[text=texto, font=\bfseries\small] at (3.5, -0.7) {Setor Alfa - Monitoramento};
\end{tikzpicture}
\end{adjustbox}
\end{center}
Baseando-se estritamente na imagem obtida pelo sensor do Setor Alfa, determine a razão matemática (em formato de fração irredutível) entre o número de veículos vermelhos e o número total de veículos estacionados nesse setor.
\par\vspace{\espacoQuestao}

\noindent\textcolor{CorLinha}{\rule{\linewidth}{0.5pt}} 
\par\vspace{\espacoPreQuestao}
\textbf{38.} A densidade demográfica é uma razão fundamental utilizada em geografia e planejamento urbano, calculada dividindo-se o número de habitantes de uma região pela sua área territorial. Se um município abriga $ \textbf{45.000 habitantes} $ em uma área exata de $ \textbf{1.500 km}^{2} $, determine a densidade demográfica local.
\par\vspace{\espacoQuestao}

\noindent\textcolor{CorLinha}{\rule{\linewidth}{0.5pt}} 
\par\vspace{\espacoPreQuestao}
\textbf{39.} (Estilo VUNESP) O cálculo da velocidade média de um veículo é uma razão matemática que relaciona a distância percorrida pelo tempo gasto no trajeto. Em uma viagem, uma família percorreu exatos $ \textbf{240 quilômetros} $ de forma ininterrupta durante $ \textbf{3 horas} $. Qual foi a razão, em $ \textbf{km/h} $, que representou a velocidade média desse automóvel?
\par\vspace{\espacoQuestao}

\noindent\textcolor{CorLinha}{\rule{\linewidth}{0.5pt}} 
\par\vspace{\espacoPreQuestao}
\textbf{40.} Um cartógrafo desenhou o mapa turístico de uma trilha utilizando regras matemáticas de proporcionalidade para garantir a fidelidade do terreno real no documento de papel.
\begin{center}
% ==========================================
% CONTROLE INDIVIDUAL DESTA IMAGEM:
% 1. CORTAR BORDAS: Mude os zeros do trim (Esquerda, Baixo, Direita, Topo). Ex: trim=1cm 0cm 1cm 0cm
% 2. TAMANHO: Para encolher só esta imagem, troque \larguraTikz por um valor (Ex: 0.5\linewidth)
% ==========================================
\begin{adjustbox}{trim=0cm 0cm 0cm 0cm, clip, max width=\larguraTikz}
\begin{tikzpicture}
\definecolor{mapa}{HTML}{FDFEFE}
\definecolor{trilha}{HTML}{27AE60}
\definecolor{moldura}{HTML}{2C3E50}
\definecolor{regua}{HTML}{F1C40F}

% --- APLICANDO DROP SHADOW ---
% LAYER 1: Base do Mapa
\fill[mapa, rounded corners=3pt, drop shadow={opacity=0.3, shadow xshift=2pt, shadow yshift=-2pt}] (0,0) rectangle (8,4);
\draw[moldura, thick, rounded corners=3pt] (0,0) rectangle (8,4);

% LAYER 2: Trilha
\draw[dashed, line width=1.5pt, trilha] (1.5, 2) -- (5.5, 2);
\fill[moldura] (1.5, 2) circle (3pt) node[above left, font=\bfseries] {Base};
\fill[moldura] (5.5, 2) circle (3pt) node[above right, font=\bfseries] {Cachoeira};

% LAYER 3: Régua na planta
\draw[|<->|, >=latex, thick] (1.5, 1.5) -- (5.5, 1.5) node[midway, fill=mapa, inner sep=2pt, font=\bfseries] {Distância Gráfica = 8 cm};

% LAYER 4: Selo de Escala
\fill[moldura, rounded corners=2pt] (0.2, 0.2) rectangle (3.2, 0.8);
\node[text=white, font=\bfseries] at (1.7, 0.5) {Escala 1 : 50.000};
\end{tikzpicture}
\end{adjustbox}
\end{center}
Considerando que a razão estipulada pela escala é de $ \textbf{1 : 50.000} $ (isto é, $ 1 $ cm no mapa equivale a $ 50.000 $ cm na realidade), calcule qual é a distância real percorrida na trilha entre a Base e a Cachoeira, entregando o resultado convertido em quilômetros.
\par\vspace{\espacoQuestao}
\noindent\textcolor{CorLinha}{\rule{\linewidth}{0.5pt}} 
\par\vspace{\espacoPreQuestao}
\textbf{41.} Para calcular a razão matemática entre duas grandezas, é obrigatório que ambas estejam expressas na mesma unidade de medida. Sabendo disso, determine a razão, na forma de fração irredutível, entre a massa de $ \textbf{450 g} $ de açúcar e a massa de um pacote de farinha que pesa $ \textbf{1,5 kg} $.
\par\vspace{\espacoQuestao}

\noindent\textcolor{CorLinha}{\rule{\linewidth}{0.5pt}} 
\par\vspace{\espacoPreQuestao}
\textbf{42.} (Estilo VUNESP) O consumo de combustível de um automóvel é diretamente proporcional à distância percorrida. O manual do proprietário informa que o veículo consome $ \textbf{8 litros} $ de gasolina para percorrer $ \textbf{100 km} $ em rodovias. Mantendo esse mesmo padrão de desempenho, determine a quantidade exata de combustível, em litros, necessária para concluir uma viagem de $ \textbf{250 km} $.
\par\vspace{\espacoQuestao}

\noindent\textcolor{CorLinha}{\rule{\linewidth}{0.5pt}} 
\par\vspace{\espacoPreQuestao}
\textbf{43.} A planta baixa a seguir apresenta a área destinada à construção de uma piscina retangular em um clube recreativo. O arquiteto desenhou o projeto utilizando instrumentos precisos de rascunho geométrico e aplicou uma escala de redução de $ 100 $ vezes.
\begin{center}
% ==========================================
% CONTROLE INDIVIDUAL DESTA IMAGEM:
% 1. CORTAR BORDAS: Mude os zeros do trim (Esquerda, Baixo, Direita, Topo). Ex: trim=1cm 0cm 1cm 0cm
% 2. TAMANHO: Para encolher só esta imagem, troque \larguraTikz por um valor (Ex: 0.5\linewidth)
% ==========================================
\begin{adjustbox}{trim=0cm 0cm 0cm 0cm, clip, max width=\larguraTikz}
\begin{tikzpicture}
\definecolor{papel}{HTML}{FDFEFE}
\definecolor{agua}{HTML}{3498DB}
\definecolor{borda}{HTML}{2C3E50}
\definecolor{cota}{HTML}{7F8C8D}

% --- APLICANDO DROP SHADOW ---
% LAYER 1: Papel do Projeto
\fill[papel, rounded corners=3pt, drop shadow={opacity=0.3, shadow xshift=2pt, shadow yshift=-2pt}] (-1, -1) rectangle (8, 4);

% LAYER 2: Piscina
\fill[agua, rounded corners=2pt, drop shadow={opacity=0.3, shadow xshift=2pt, shadow yshift=-2pt}] (1, 0.5) rectangle (6, 2.5);
\draw[borda, thick, rounded corners=2pt] (1, 0.5) rectangle (6, 2.5);

% LAYER 3: Cotas do Desenho (Linhas de medida do papel)
\draw[|<->|, >=latex, thick, cota] (1, 0) -- (6, 0) node[midway, fill=papel, inner sep=2pt, font=\bfseries] {12 cm no desenho};
\draw[|<->|, >=latex, thick, cota] (6.5, 0.5) -- (6.5, 2.5) node[midway, fill=papel, inner sep=2pt, font=\bfseries, rotate=90] {5 cm};

% LAYER 4: Escala Indicada
\node[fill=borda, text=white, rounded corners=2pt, inner sep=3pt, font=\bfseries\small] at (1.5, 3.2) {Escala 1 : 100};
\end{tikzpicture}
\end{adjustbox}
\end{center}
Determine as dimensões reais da piscina (comprimento e largura) em metros e, em seguida, calcule a sua área real exata em metros quadrados ($ \textbf{m}^{2} $).
\par\vspace{\espacoQuestao}

\noindent\textcolor{CorLinha}{\rule{\linewidth}{0.5pt}} 
\par\vspace{\espacoPreQuestao}
\textbf{44.} (ENEM - Adaptada) A velocidade média é a grandeza que expressa a razão entre a distância total percorrida e o tempo total de viagem. Um ônibus de viagem percorreu a distância intermunicipal de $ \textbf{120 km} $ de forma constante durante um intervalo de exatas $ \textbf{1,5 hora} $. Assinale a alternativa que indica a velocidade média desenvolvida pelo veículo. \textbf{Em seguida, justifique o cálculo, refutando o porquê de um aluno desatento possivelmente assinalar a alternativa (a).}
\begin{enumerate}[label=\alph*)]
    \item $ 180 \text{ km/h} $
    \item $ 60 \text{ km/h} $
    \item $ 80 \text{ km/h} $
    \item $ 100 \text{ km/h} $
\end{enumerate}
\par\vspace{\espacoQuestao}

\noindent\textcolor{CorLinha}{\rule{\linewidth}{0.5pt}} 
\par\vspace{\espacoPreQuestao}
\textbf{45.} A propriedade fundamental das proporções afirma que o produto dos meios é igual ao produto dos extremos. Aplique essa propriedade rigorosamente para resolver a proporção matemática a seguir, determinando o valor numérico exato da variável $ \textbf{x} $.
$$ \mfrac{\textbf{x}}{15} = \mfrac{4}{5} $$
\par\vspace{\espacoQuestao}

\noindent\textcolor{CorLinha}{\rule{\linewidth}{0.5pt}} 
\par\vspace{\espacoPreQuestao}
\textbf{46.} A interface gráfica de um aplicativo de transferência de arquivos exibe o progresso de um download diretamente no painel do usuário.
\begin{center}
% ==========================================
% CONTROLE INDIVIDUAL DESTA IMAGEM:
% 1. CORTAR BORDAS: Mude os zeros do trim (Esquerda, Baixo, Direita, Topo). Ex: trim=1cm 0cm 1cm 0cm
% 2. TAMANHO: Para encolher só esta imagem, troque \larguraTikz por um valor (Ex: 0.5\linewidth)
% ==========================================
\begin{adjustbox}{trim=0cm 0cm 0cm 0cm, clip, max width=\larguraTikz}
\begin{tikzpicture}
\definecolor{tela}{HTML}{1A1A1A}
\definecolor{barravazia}{HTML}{7F8C8D}
\definecolor{barracheia}{HTML}{2ECC71}
\definecolor{texto}{HTML}{ECF0F1}

% --- APLICANDO DROP SHADOW ---
% LAYER 1: Moldura do Aplicativo
\fill[tela, rounded corners=4pt, drop shadow={opacity=0.4, shadow xshift=3pt, shadow yshift=-3pt}] (0,0) rectangle (7, 2.5);

% LAYER 2: Barra de Progresso
\fill[barravazia, rounded corners=2pt] (1, 1.2) rectangle (6, 1.6);
\fill[barracheia, rounded corners=2pt] (1, 1.2) rectangle (3, 1.6); % 40% da barra

% LAYER 3: Informações de Texto
\node[text=texto, font=\bfseries] at (3.5, 2) {Progresso do Download};
\node[text=texto, font=\bfseries] at (1.5, 0.7) {40\% Concluído};
\node[text=barracheia, font=\bfseries] at (5.2, 0.7) {Baixado: 120 MB};
\end{tikzpicture}
\end{adjustbox}
\end{center}
Baseando-se nos dados numéricos exibidos no painel do aplicativo, utilize o raciocínio proporcional para calcular o tamanho total exato desse arquivo em Megabytes ($ \textbf{MB} $).
\par\vspace{\espacoQuestao}

\noindent\textcolor{CorLinha}{\rule{\linewidth}{0.5pt}} 
\par\vspace{\espacoPreQuestao}
\textbf{47.} Para preparar a coloração de uma fachada, um pintor mistura tinta acrílica azul e tinta acrílica amarela na razão estrita de $ \textbf{2 partes} $ de azul para cada $ \textbf{5 partes} $ de amarelo. Se ele despejar no balde de mistura exatos $ \textbf{750 ml} $ de tinta amarela, determine a quantidade proporcional de tinta azul que deverá ser adicionada para manter a tonalidade original.
\par\vspace{\espacoQuestao}

\noindent\textcolor{CorLinha}{\rule{\linewidth}{0.5pt}} 
\par\vspace{\espacoPreQuestao}
\textbf{48.} O aluno fictício Thiago precisava calcular a densidade demográfica de um pequeno município que possui $ \textbf{5.000 habitantes} $ e uma área territorial de $ \textbf{250 km}^{2} $. Ele escreveu o seguinte cálculo na lousa:
\textit{"Densidade é uma divisão. Se eu dividir a área pelo número de pessoas, terei $ 250 \div 5000 = 0,05 $."}
\begin{boxresponda}
\textbf{Responda:} Thiago errou a modelagem conceitual da razão. Explique textualmente o que significa densidade demográfica, identifique o erro na montagem da fração dele e efetue o cálculo correto apresentando a unidade de medida adequada.
\end{boxresponda}
\par\vspace{\espacoQuestao}

\noindent\textcolor{CorLinha}{\rule{\linewidth}{0.5pt}} 
\par\vspace{\espacoPreQuestao}
\textbf{49.} Um barista prepara uma bebida fria artesanal misturando concentrado de café expresso e leite evaporado. Os recipientes medidores abaixo ilustram a dosagem exata extraída para a montagem de um único copo da bebida.
\begin{center}
% ==========================================
% CONTROLE INDIVIDUAL DESTA IMAGEM:
% 1. CORTAR BORDAS: Mude os zeros do trim (Esquerda, Baixo, Direita, Topo). Ex: trim=1cm 0cm 1cm 0cm
% 2. TAMANHO: Para encolher só esta imagem, troque \larguraTikz por um valor (Ex: 0.5\linewidth)
% ==========================================
\begin{adjustbox}{trim=0cm 0cm 0cm 0cm, clip, max width=\larguraTikz}
\begin{tikzpicture}
\definecolor{fundo}{HTML}{F8F9F9}
\definecolor{cafe}{HTML}{5D4037}
\definecolor{leite}{HTML}{e7e7e7}
\definecolor{vidro}{HTML}{BDC3C7}
\definecolor{texto}{HTML}{2C3E50}

% --- APLICANDO DROP SHADOW ---
% LAYER 1: Bancada
\fill[fundo, drop shadow={opacity=0.2, shadow yshift=-2pt}] (-1, 0) rectangle (7, 5);

% LAYER 2: Béquer do Café
\fill[cafe, drop shadow={opacity=0.3, shadow xshift=2pt}, rounded corners=1pt] (0.5, 0.5) rectangle (2.5, 2);
\draw[vidro, thick, rounded corners=2pt] (0.5, 0.5) rectangle (2.5, 3.5);
\foreach \y/\val in {1.5/100, 2/150, 2.5/200, 3/250} {
    \draw[texto, thick] (0.5, \y) -- (0.8, \y);
}
\node[text=texto, font=\bfseries\small] at (1.5, 2.3) {150 ml};
\node[text=texto, font=\bfseries] at (1.5, -0.3) {Extrato de Café};

% LAYER 3: Béquer do Leite
\fill[leite, opacity=0.9, drop shadow={opacity=0.3, shadow xshift=2pt}, rounded corners=1pt] (4, 0.5) rectangle (6, 3);
\draw[vidro, thick, rounded corners=2pt] (4, 0.5) rectangle (6, 4);
\foreach \y/\val in {1.5/100, 2.5/200, 3/250, 3.5/300} {
    \draw[texto, thick] (4, \y) -- (4.3, \y);
}
\node[text=texto, font=\bfseries\small] at (5, 3.3) {250 ml};
\node[text=texto, font=\bfseries] at (5, -0.3) {Leite Evaporado};
\end{tikzpicture}
\end{adjustbox}
\end{center}
A partir dos volumes extraídos visualmente dos medidores, estabeleça a razão matemática entre a quantidade de extrato de café e a quantidade de leite evaporado. Entregue a resposta na sua forma irredutível.
\par\vspace{\espacoQuestao}

\noindent\textcolor{CorLinha}{\rule{\linewidth}{0.5pt}} 
\par\vspace{\espacoPreQuestao}
\textbf{50.} Organize as ideias: na matemática, dizemos que duas frações formam uma proporção verdadeira quando elas são equivalentes. Explique de forma prática, utilizando a ideia de simplificação de frações, por que a razão $ \mfrac{4}{6} $ representa exatamente a mesma proporção que $ \mfrac{10}{15} $.
\par\vspace{\espacoQuestao}

\noindent\textcolor{CorLinha}{\rule{\linewidth}{0.5pt}} 
\par\vspace{\espacoPreQuestao}
\textbf{51.} (ENEM - Adaptada) A cozinheira de um restaurante industrial segue uma ficha técnica que indica o uso de $ \textbf{250 g} $ de farinha de rosca para empanar filés suficientes para atender uma mesa com $ \textbf{4 pessoas} $. Certo dia, ela recebeu uma reserva exclusiva para um banquete com $ \textbf{10 pessoas} $. Assinale a alternativa que traz a quantidade proporcional exata de farinha de rosca que ela deverá utilizar.
\begin{enumerate}[label=\alph*)]
    \item $ 400 \text{ g} $
    \item $ 500 \text{ g} $
    \item $ 550 \text{ g} $
    \item $ 625 \text{ g} $
\end{enumerate}
\par\vspace{\espacoQuestao}

\noindent\textcolor{CorLinha}{\rule{\linewidth}{0.5pt}} 
\par\vspace{\espacoPreQuestao}
\textbf{52.} O computador de bordo de um veículo de frota empresarial armazena os dados das viagens realizadas para auditar a eficiência energética do motorista. A tela exibe o relatório da última entrega intermunicipal concluída.
\begin{center}
% ==========================================
% CONTROLE INDIVIDUAL DESTA IMAGEM:
% 1. CORTAR BORDAS: Mude os zeros do trim (Esquerda, Baixo, Direita, Topo). Ex: trim=1cm 0cm 1cm 0cm
% 2. TAMANHO: Para encolher só esta imagem, troque \larguraTikz por um valor (Ex: 0.5\linewidth)
% ==========================================
\begin{adjustbox}{trim=0cm 0cm 0cm 0cm, clip, max width=\larguraTikz}
\begin{tikzpicture}
\definecolor{painel}{HTML}{2C3E50}
\definecolor{visor}{HTML}{ECF0F1}
\definecolor{destaque}{HTML}{3498DB}
\definecolor{texto}{HTML}{2C3E50}

% --- APLICANDO DROP SHADOW ---
% LAYER 1: Estrutura do Painel Veicular
\fill[painel, rounded corners=6pt, drop shadow={opacity=0.4, shadow xshift=3pt, shadow yshift=-3pt}] (0,0) rectangle (9, 3.5);

% LAYER 2: Tela Digital Principal
\fill[visor, rounded corners=3pt, drop shadow={opacity=0.2}] (0.3, 0.3) rectangle (8.7, 3.2);

% LAYER 3: Divisórias e Informações
\draw[painel, thick, dashed] (4.5, 0.5) -- (4.5, 3);
\node[text=texto, font=\bfseries\Large] at (2.25, 2) {360 km};
\node[text=texto, font=\bfseries\small] at (2.25, 1.2) {Distância Percorrida};

\node[text=destaque, font=\bfseries\Large] at (6.5, 2) {30 Litros};
\node[text=texto, font=\bfseries\small] at (6.5, 1.2) {Combustível Consumido};

\node[fill=painel, text=white, rounded corners=2pt, inner sep=2pt, font=\bfseries\footnotesize] at (4.5, 2.9) {RELATÓRIO DE ROTA};
\end{tikzpicture}
\end{adjustbox}
\end{center}
A eficiência energética (ou consumo médio) de um carro no Brasil é mensurada pela razão entre a distância percorrida e a quantidade de combustível consumido. Utilizando os dados da tela digital, calcule a eficiência dessa viagem em quilômetros por litro ($ \textbf{km/L} $).

\end{multicols*}
\end{document}
